\section{A Tiered Detection Architecture}
\label{sec:architecture}

\Cref{prop:impossible,prop:tiers} together prescribe the architecture:
governance-drift detection is configuration comparison \emph{plus} the
minimal additional evidence streams, joined against the selected admitted basis.
\Cref{fig:architecture} organizes it as cumulative tiers, each named for what
it adds and each mapped to the drift classes it first decides.

\begin{figure*}[t]
\centering
\resizebox{0.93\textwidth}{!}{% Tiered detection architecture (single column, resized)
\begin{tikzpicture}[node distance=2.6mm and 4mm]
  \node[attn, minimum width=25mm] (snap) at (0,0)
    {\textbf{append-only}\\\textbf{approval recorder}\\{\smlbl $G^{(0)}\to\cdots\to B_x(t)$}\\{\smlbl $m_B,\delta_B,\pi_B,A_B^{\rm proof},\iota_B,\sigma_B$}};

  % tiers stack
  \node[stage, right=9mm of snap, minimum width=26mm, yshift=13mm] (t0)
    {T0/T0n: config comparison\\{\smlbl reconciler loop \; $\to$ S1}};
  \node[stage, below=2mm of t0, minimum width=26mm] (t1)
    {T1: $+$ policy stream $\Pol(t)$\\{\smlbl policy repo \; $\to$ S3}};
  \node[stage, below=2mm of t1, minimum width=26mm] (t2)
    {T2: $+$ authorization\\{\smlbl proof + $\Auth_{\rm live}(t)$ \; $\to$ S2, S6(intent), S9}};
  \node[stage, below=2mm of t2, minimum width=26mm] (t3)
    {T3: $+$ artifact lineage\\{\smlbl digests, attestations \; $\to$ S4, S6(authz), S8}};
  \node[stage, below=2mm of t3, minimum width=26mm] (t4)
    {T4: $+$ environment inventory\\{\smlbl cloud audit/inventory \; $\to$ S5, S7}};

  % evaluator
  \node[evid, right=9mm of t2, minimum width=22mm] (eval)
    {\textbf{evaluator}\\{\smlbl computes $(\GDsub,\ObsMask)$}\\{\smlbl per deployment}};
  \node[brok, below=3mm of eval, minimum width=22mm] (out)
    {full vector + priority\\{\smlbl component $\to$ remediation}\\{\smlbl $\unk$ where evidence fails}};

  \draw[eflow] (snap.east) to[out=0,in=180] (eval.west);
  \foreach \n in {t0,t1,t2,t3,t4}
    \draw[flow] (\n.east) to[out=0,in=180] (eval.west);
  \draw[flow] (eval) -- (out);
\end{tikzpicture}
}
\caption{Tiered governance-drift detection. The append-only recorder (left)
preserves successive approvals and selects $B_x(t)$. Evidence tiers (center) cumulatively widen
the detector's vocabulary: T0/T0n configuration comparison (the reconciler's
existing loop, naive or churn-normalized), T1 policy-version stream, T2
authorization records with validity semantics, T3 artifact lineage, T4
environment inventory. The evaluator (right) recomputes the component
divergences of \cref{sec:model-distance} continuously and emits classed drift
vectors plus an epistemic mask, with a priority projection and the remediation
family each component implies. Stream adapters also report the presence,
completeness, freshness, integrity, authenticity, subject linkage, and health
required by each component's evidence contract. Annotations show which components of
Scenarios S1--S9 each tier first detects, as measured in
\cref{sec:study}.}
\label{fig:architecture}
\end{figure*}

\subsection{Components}
\textbf{Approved-state recorder.} At approval, append an immutable pending
snapshot to $\Gapp^x$---manifest content, \emph{resolved} digests, policy-version pin,
immutable authorization proof with execution times, intent linkage, environment
assumptions, scope, and predecessor edge. At operational effect, append or
attest its activation; cancellation marks it aborted. Persist both events
durably and independently of the delivery
platforms. This is the architecture's one hard prerequisite: without a
recorded approval basis there is nothing to drift \emph{from}; where such
records are absent or decayed, the basis-dependent entries of
$\ObsMask_{\mu}(u,\theta;T)$ are zero and those components are undecidable. If
no independently evaluated component is known nonzero, the total status is
undecidable; if configuration or another independently established component
is known nonzero, the total status is inconsistent because known inconsistency
precedes unknown evidence. The desired/observed configuration component remains
comparable when its own evidence contract is satisfied; the evaluator must
report this partial vector and mask rather than either total consistency or
total ignorance.
Conflicting activated maxima also fail closed instead of being
ordered by approval timestamp (\cref{prop:basis-selection}).

\textbf{Evidence tiers.}
\emph{T0/T0n}: the desired/observed comparison every reconciler already
performs~\cite{argocd,fluxcd,awsconfig}; T0n normalizes runtime-owned fields
(replica counts under autoscaling, rollout hashes)---the study quantifies why
the naive form is unusable as a governance signal.
\emph{T1}: the policy-version stream from the policy engine's
repository~\cite{opa,kyverno}; the evaluator re-evaluates the running state
against $\Pol(t)$ and compares with the pinned basis.
\emph{T2}: authorization records with mode and validity semantics---expiry,
scope, prospective or retroactive revocation---bound to the expected change
graph, covered digests where the authorization scope requires them, and
immutable approval/exception proofs. Exceptions and approvals are therefore
evidenced objects with clocks, subjects, and coverage rather than booleans;
ITIL-style change records~\cite{itil4} and platform approvals both map here.
\emph{T3}: artifact lineage---digest-level linkage from running workloads to
approval subjects, the slice pioneered by supply-chain
attestation~\cite{slsa,torres2019,newman2022,binauthz}.
\emph{T4}: environment inventory---observation of the unmanaged properties
assumed at approval (IAM scope, exposure) and comparison against the
\emph{recorded} predicates $\sigma_B$ in $B_x(t)$, from cloud inventories,
audit streams, and telemetry correlation~\cite{cloudtrail,cspm,otel}.
Mechanically this is baseline comparison---CSPM's move---and we say so
plainly; the difference is the reference (the approval event's recorded
assumptions rather than a benchmark) and the joined, class-resolved verdict.

\textbf{Evaluator.} A continuous process (or scheduled sweep) computing
$\GD_{\mu}(u,\theta;T)$ per deployment: cheap incremental checks on context-change events
(policy revision published, exception clock elapsed, registry event), using
event-triggered recomputation where a source exposes deltas and scheduled or
background re-evaluation where it does not. Several scenarios
in \cref{sec:study} are detectable with zero latency at the tier that
carries the relevant modeled event; this is not a claim that live integrations
avoid polling or scans.

\textbf{Analytic scaling envelope.} For $n$ deployments with at most $m$
managed configuration fields, $r$ policy predicates, $a$ relied-upon
instruments, $d$ running digests, and $q$ environment predicates each, a full
sweep is $O(n(m+r+a+d+q))$ time. Retaining $s$ complete snapshots per
deployment costs $O(ns(m+a+d+q))$ storage, plus current-stream indices. The
intended execution is incremental: policy-to-deployment,
subject-to-approval, digest-to-workload, and predicate-to-resource indices
restrict an event to its $k$ affected deployments, yielding
$O(k(m+r+a+d+q))$ work. Policy results and immutable proofs are cacheable;
evaluators can shard by cluster, namespace, or tenant. Kubernetes watches and
batched cloud-inventory deltas avoid one API request per deployment. Cloud
quotas and cache invalidation remain engineering constraints. The in-memory
microbenchmark in \cref{tab:scaling} measures the join after decoding and
indexing inputs; it deliberately does not claim a Kubernetes or production
throughput bound.

\textbf{Minimal adoption path.} The smallest useful deployment is
T0n+T1+T2: normalize the reconciler comparison, re-evaluate running state
under current policy, and join durable approvals and exceptions with explicit
validity semantics. This separates operational, compliance, and authorization
risk immediately; T3 adds digest lineage and T4 adds declared environment
predicates as evidence maturity grows. Governance Conformance Coverage (GCC;
\cref{met:gcc}) is reported from the first rollout so unevaluable coverage is
visible rather than treated as conformance.

For legacy workloads, a discovery pass may create a signed
\emph{retrospective baseline} from current state, current policy results, and
an accountable owner's attestation. It is timestamped and labeled as
reconstructed evidence---never represented as a historical approval---and
GCC remains partial for facts that cannot be recovered. New changes then enter
the ordinary append-only path, allowing coverage to increase without
rewriting the estate's past.

\subsection{Threat and Trust Model}
The protected claim is integrity and decidability of the relation, not
availability of the workload. An attacker may mutate Git, tags, runtime
objects, policy results, authorization status, inventory, or one evidence
adapter; operators may also make the same changes accidentally. These sources
are therefore evidence inputs, not intrinsically trusted truth. The trusted
computing base (TCB) is the recorder/activation authority, stable-identity
issuer, authenticated time and adapter metadata, append-only log or external
anchor, basis selector, and evaluator code. Every record is subject-bound,
sequenced, freshness-bounded, and content-addressed; snapshots are complete
atomic tuples chained to their predecessors. A contract failure masks only
the dependent components and prevents a conformance claim.

A local hash chain detects modification and reordering of retained snapshots,
but cannot by itself prove that a privileged actor did not truncate the tail
or replace the entire history. Production deployment therefore requires an
external transparency anchor or threshold signatures, access separation, and
key/clock rotation records. Compromise of the evaluator and all anchors, or
collusion of the recorder and activation authority, is outside the model.
Denial of evidence remains visible as reduced GCC/undecidable; remediation may
keep a critical workload available, but it must not relabel uncertainty as
conformance (\cref{sec:discussion}).

\subsection{Relation to Existing Mechanisms}
The architecture is deliberately assembled from streams that already exist in
mature platforms; its novelty is the \emph{join semantics} against a recorded
approval basis and the class-resolved verdict. Admission
control~\cite{k8sadmission,opa,kyverno} evaluates the same predicates once,
at entry; continuous-validation features~\cite{binauthz} re-evaluate the
artifact-policy slice; posture management~\cite{cspm} scans environment
configuration against benchmark policy without any approved-state reference.
Indeed, per-tier continuous detection largely \emph{exists} in cited
tooling, and the architecture credits it explicitly: Kyverno's background
scanning re-evaluates existing resources against current policy (our
T1)~\cite{kyverno}; configuration recorders track unmanaged resources
(much of T4)~\cite{awsconfig,cspm}; continuous validation re-checks
artifact policy (part of T3)~\cite{binauthz}. Among the mechanisms in the
bounded source set documented in \cref{sec:related}, none jointly provides all
of the following: every tier evaluated against the
\emph{admitted basis} $B_x(t)$, with mode-aware validity on the authorization
itself, and verdicts that distinguish ``violates policy'' from ``was never
approved'' from ``was admitted legitimately and the world moved''---the
distinctions that determine remediation (\cref{sec:model-vector}). The
contribution is the relation and its join semantics, not the comparison
machinery, and \cref{sec:related} restates novelty in exactly those terms.
